% Options for packages loaded elsewhere
\PassOptionsToPackage{unicode}{hyperref}
\PassOptionsToPackage{hyphens}{url}
\documentclass[
]{article}
\usepackage{xcolor}
\usepackage[margin=1in]{geometry}
\usepackage{amsmath,amssymb}
\setcounter{secnumdepth}{-\maxdimen} % remove section numbering
\usepackage{iftex}
\ifPDFTeX
  \usepackage[T1]{fontenc}
  \usepackage[utf8]{inputenc}
  \usepackage{textcomp} % provide euro and other symbols
\else % if luatex or xetex
  \usepackage{unicode-math} % this also loads fontspec
  \defaultfontfeatures{Scale=MatchLowercase}
  \defaultfontfeatures[\rmfamily]{Ligatures=TeX,Scale=1}
\fi
\usepackage{lmodern}
\ifPDFTeX\else
  % xetex/luatex font selection
\fi
% Use upquote if available, for straight quotes in verbatim environments
\IfFileExists{upquote.sty}{\usepackage{upquote}}{}
\IfFileExists{microtype.sty}{% use microtype if available
  \usepackage[]{microtype}
  \UseMicrotypeSet[protrusion]{basicmath} % disable protrusion for tt fonts
}{}
\makeatletter
\@ifundefined{KOMAClassName}{% if non-KOMA class
  \IfFileExists{parskip.sty}{%
    \usepackage{parskip}
  }{% else
    \setlength{\parindent}{0pt}
    \setlength{\parskip}{6pt plus 2pt minus 1pt}}
}{% if KOMA class
  \KOMAoptions{parskip=half}}
\makeatother
\usepackage{graphicx}
\makeatletter
\newsavebox\pandoc@box
\newcommand*\pandocbounded[1]{% scales image to fit in text height/width
  \sbox\pandoc@box{#1}%
  \Gscale@div\@tempa{\textheight}{\dimexpr\ht\pandoc@box+\dp\pandoc@box\relax}%
  \Gscale@div\@tempb{\linewidth}{\wd\pandoc@box}%
  \ifdim\@tempb\p@<\@tempa\p@\let\@tempa\@tempb\fi% select the smaller of both
  \ifdim\@tempa\p@<\p@\scalebox{\@tempa}{\usebox\pandoc@box}%
  \else\usebox{\pandoc@box}%
  \fi%
}
% Set default figure placement to htbp
\def\fps@figure{htbp}
\makeatother
\setlength{\emergencystretch}{3em} % prevent overfull lines
\providecommand{\tightlist}{%
  \setlength{\itemsep}{0pt}\setlength{\parskip}{0pt}}
\usepackage{booktabs}
\usepackage{longtable}
\usepackage{array}
\usepackage{multirow}
\usepackage{wrapfig}
\usepackage{float}
\usepackage{colortbl}
\usepackage{pdflscape}
\usepackage{tabu}
\usepackage{threeparttable}
\usepackage{threeparttablex}
\usepackage[normalem]{ulem}
\usepackage{makecell}
\usepackage{xcolor}
\usepackage{bookmark}
\IfFileExists{xurl.sty}{\usepackage{xurl}}{} % add URL line breaks if available
\urlstyle{same}
\hypersetup{
  pdftitle={baby\_names\_pdf},
  hidelinks,
  pdfcreator={LaTeX via pandoc}}

\title{baby\_names\_pdf}
\author{}
\date{\vspace{-2.5em}2026-06-19}

\begin{document}
\maketitle

\newpage

\section{Introduction}\label{introduction}

This report analyzes historical trends, pop-culture shocks, and media
genre influences within national baby-naming datasets to inform
long-term inventory forecasting and intellectual property licensing
strategies. F By tracking how fast name trends persist, peak, and decay
across different entertainment channels, this analysis establishes a
data-driven framework for timing product lifecycles and optimizing
personalized toy lines.

\section{Persistence analysis}\label{persistence-analysis}

\begin{verbatim}
## $timeseries
## # A tibble: 618 x 4
##     year gender   lag   rho
##    <dbl> <chr>  <int> <dbl>
##  1  1910 F          1 0.979
##  2  1910 F          2 0.951
##  3  1910 F          3 0.942
##  4  1910 M          1 0.975
##  5  1910 M          2 0.922
##  6  1910 M          3 0.925
##  7  1911 F          1 0.975
##  8  1911 F          2 0.933
##  9  1911 F          3 0.915
## 10  1911 M          1 0.955
## # i 608 more rows
## 
## $plot
\end{verbatim}

\pandocbounded{\includegraphics[keepaspectratio]{baby_names_pdf_files/figure-latex/unnamed-chunk-1-1.pdf}}

\begin{verbatim}
## 
## $era_comparison
## # A tibble: 6 x 5
##   gender   lag from_1990 pre_1990 drop_in_persistence
##   <chr>  <int>     <dbl>    <dbl>               <dbl>
## 1 F          1     0.937    0.952              0.0143
## 2 F          2     0.868    0.900              0.0318
## 3 F          3     0.815    0.853              0.0379
## 4 M          1     0.950    0.981              0.0306
## 5 M          2     0.891    0.958              0.0669
## 6 M          3     0.836    0.934              0.0981
\end{verbatim}

The chart shows that girls' names dropped much faster in persistence
post-1990 than boys' names. A toy brand relying on the same names for
girls' products year-over-year faces a higher risk of sounding outdated.
\# Popularity surges

This plot shows the names for which the year-on-year growth has doubled.
An observation is coloured red if the name doubled in a two range of the
release rate of a tv show or movie and in coloured blue if the name
doubled in the a two range since a artist reached there best position on
the Top 100. This clearly illustrates how quickly pop culture affects
Baby\_Names. Looking at the left-bottom we see Zendaya is name that has
grown and has played in many movies recently.

\section{Density}\label{density}

\pandocbounded{\includegraphics[keepaspectratio]{baby_names_pdf_files/figure-latex/unnamed-chunk-3-1.pdf}}
The ``Flash-in-the-Pan'' Risk (Arya, Khaleesi, Miley): These names show
sharp, vertical spikes with almost zero historical runway. They
represent intense, immediate demand that risks leaving you with dead
inventory if you over-manufacture customized toys past the peak of the
media hype.

The ``Sustained Icon'' Value (Billie, Drake, Rihanna): The red dashed
lines highlight that these names were already rising or heavily
established before their massive pop-culture breakout. These trends have
a smoother decay curve, making them safer, longer-term investments for
personalized products.

The ``Leading Indicator'' Opportunity (Zendaya): The name spikes exactly
at the red line, signaling a highly concentrated generation of children
with this name hitting your target toddler-to-youth toy market
simultaneously.

\section{Genres}\label{genres}

\pandocbounded{\includegraphics[keepaspectratio]{baby_names_pdf_files/figure-latex/unnamed-chunk-4-1.pdf}}

High-Yield Toy Lines (Drama \& Comedy): Shows categorized under Drama
(27 spikes) and Comedy (20 spikes) are your most fertile ground for
long-term IP licensing. Characters in these genres successfully capture
deep consumer attachment, driving parents to name their children after
them.

Low-Yield Genre Risks (Family, Sci-Fi, Fantasy): Surprisingly,
traditional toy-driving genres like Family (2) and Fantasy (7) trigger
fewer historical baby name spikes here. This suggests that while these
genres sell toys immediately, they don't always inspire the lifelong
emotional investment required to influence a child's legal name.

\section{Conclusion}\label{conclusion}

We see that pop culture has a massive influence on baby names. Very
popular name these days are Arya, Khaleesi(from Game of thrones) or
Zendaya. It is also noted that name persistence has decreased post
1990s. This means that toy manufactures should not sign as big contracts
relative to population size compared to before the 1990s

\end{document}
